\section{Turn 1: the complete joint arithmetic object}
This tranche addresses simultaneous exterior/middle failure in every original
first-half shell. It places no chamber, Farey-depth, residual-primality, or
dyadic-seed restriction on that shell. The output is a structural reduction,
not a theorem that every prime has an occupied shell.

The prime-shift three-target budget and primitive index-six exceptional factor
already occur in the centl notes \cite[\S\S4--9]{BryanTwo},
\cite[\S\S3--8]{BryanSix}. They are not claimed as new here. We extend the
calculation to arbitrary residual moduli, distinguish the full and effective
quotients, pull the targets through every factor cut, classify all effective
quotients of order at most eight, and retain the exact kernel coefficients
and extension cocycle. Every assertion used below has a proof in this tranche.

Fix a prime $p\equiv1\pmod4$ and an integer $p/4<a<p/2$. Put
\[
 R=4a-p,\qquad a=\prod_{i=1}^r q_i^{e_i},\qquad
 \mathcal B=\prod_i[-e_i,e_i]_{\mathbb Z},\qquad
 v(\beta)=\prod_iq_i^{\beta_i}\pmod R.                 \tag{A1}
\]
The $q_i$ are distinct actual primes. Since $R\equiv3\pmod4$ and
$\gcd(pa,R)=1$, all inverses in (A1) exist. Retain
\[
 u(\beta)=\prod_iq_i^{e_i+\beta_i},\quad
 \mu(g)=\#\{\beta\in\mathcal B:v(\beta)=g\},\quad W=\operatorname{supp}\mu.
                                                               \tag{A2}
\]
The inverse is $\beta_i=v_{q_i}(u)-e_i$. In particular
$\sum_g\mu(g)=\prod_i(2e_i+1)$ and $\mu(g)=\mu(g^{-1})$.
The channel labels and targets are
\[
 (M,-1),\quad(E_+,-p^{-1}),\quad(E_-,-p).                 \tag{A3}
\]
$E_-$ is the same exterior test after the recorded involution
$\beta\mapsto-\beta$, not a third canonical channel. Thus
\[
 |M_a|=\mu(-1),\qquad |E_a|=\mu(-p^{-1}),\qquad
 E_a=M_a=\varnothing\iff W\cap\{-1,-p,-p^{-1}\}=\varnothing. \tag{A4}
\]
The exponent vector is the same at every prime power dividing $R$.
A compatible unrestricted CRT residue is not substituted for a vector in
$\mathcal B$.

\paragraph{Original ordered return.}
For a hit, put $d_0=\gcd(a,u)$ and
\[
 (h,r_0,s_0)=(d_0^2/u,u/d_0,a/d_0).
\]
At $q_i$, their valuations are
$(e_i-|\beta_i|,(\beta_i)_+,(-\beta_i)_+)$. Hence they are positive
integers, $a=hr_0s_0$, $u=hr_0^2$, and $\gcd(r_0,s_0)=1$.
Since $u\equiv av(\beta)$ and $4a\equiv p\pmod R$, (A3) becomes
$R\mid u+a$ in M and $R\mid4u+1$ in E. All factors of $a$ are units
modulo $R$, so the required integer quotients are
\[
 \begin{array}{ll}
 M:&\lambda=(r_0+s_0)/R,\quad (x,y,z)=(a,phs_0\lambda,phr_0\lambda),\\
 E:&\kappa=(pr_0+s_0)/R,\quad (x,y,z)=(a,hs_0\kappa,phr_0\kappa).
 \end{array}                                                   \tag{A5}
\]
Multiplying the reciprocal sum by $phr_0s_0\lambda$ or
$phr_0s_0\kappa$ verifies $4/p$. The ordered inverses are
\[
 u=\frac{pa^2}{Ry-pa}\quad(M),\qquad
 u=\frac{a^2}{Ry-pa}\quad(E).                              \tag{A6}
\]
The M involution $u\mapsto a^2/u$ swaps $y,z$ and has no hit fixed point:
its fixed divisor $u=a$ would require $R\mid2a$, impossible. Thus $|M_a|$
is even. In E, the $E_-$ label is first negated before applying (A5).
These arithmetic identities are inherited Type I/II coordinates
\cite[\S2, Propositions 2.2 and 2.6]{ET}.

\paragraph{Completeness of the interval.}
Here is the required reduction, rather than an assumed smaller search range.
Order any solution as $x\le y\le z$. Positivity gives $x>p/4$.
If $x>p/2$, then $4/p\le1/x+2/y<2/p+2/y$, hence $y<p$.
The cleared equation forces $z=pk$. Then $(4k-1)xy=pk(x+y)$,
so $t=(4k-1)/p$ is a positive integer, $t\equiv3\pmod4$ and
$\gcd(t,k)=1$. Write $x=gA,y=gB$, $\gcd(A,B)=1$.
The equation gives $k=ABD$, $tg=D(A+B)$ and $\gcd(t,D)=1$.
The inequality $x>p/2$ gives $2DA(A-B)>-1$, so $A\ge B$.
Since $A\le B$, both equal one; then $tg=2D$ forces $t\mid2$, a
contradiction. Equality $x=p/2$ is impossible for odd $p$.

With $a=x$, the other denominators satisfy
$(Ry-pa)(Rz-pa)=p^2a^2$, with both factors positive. The exponent of $p$
in the first factor is $0,1$, or $2$. In case zero set
$u=a^2/(Ry-pa)$ and use E; in case two set $u=(Ry-pa)/p^2$ and
use E with $y,z$ exchanged; in case one set $u=pa^2/(Ry-pa)$ and
use M. Each is a divisor of $a^2$, and the congruence of the first factor
modulo $R$ proves its gate. This proves completeness with all permutations
retained. The same first-half argument is recorded in the original
workbench \cite[Theorem \texttt{bp-first-half}]{Prior}.

\section{The actual period quotient and every target collision}
Use the full unit group $G=(\mathbb Z/R\mathbb Z)^\times$ and the smaller
coefficient envelope
\[
 \Gamma=\langle-1,2,q_1,\ldots,q_r\rangle\le G,
 \qquad K=\operatorname{Stab}_G(W)=\{g:gW=W\}.             \tag{B1}
\]
The generators $2,-1$ in $\Gamma$ are coefficient units, not newly supplied
factors of $a$. All targets are in $\Gamma$ because $p\equiv4a$.
Since $1\in W$, every $g\in K$ lies in $W$. Therefore
$K\le\Gamma$ and $\operatorname{Stab}_\Gamma(W)=K$.
Let
\[
 A=\Gamma/K,\quad d=|A|,\quad S=W/K,\quad
 \alpha=-1K,\ z=2K,\ \eta=pK,\ g_i=q_iK.                 \tag{B2}
\]
The ambient index $[G:K]$ is separately retained; it can exceed $d$.
The exact arithmetic relation is
\[
 \boxed{\eta=z^2\prod_i g_i^{e_i}.}                      \tag{B3}
\]
The quotient support $S$ contains $1$, is inversion-invariant, and has
trivial stabilizer. To prove the last claim, lift any quotient period:
it would stabilize the full union $W$ of $K$-cosets and hence belong to $K$.

\begin{theorem}[Joint failure and collision law]\label{thm:collision}
For every original shell, joint failure is equivalent to
\[
 S\cap F=\varnothing,\qquad
 F=\{\alpha,\alpha\eta,\alpha\eta^{-1}\}.               \tag{B4}
\]
Under failure, $\alpha\ne1$ has order two, $d$ is even, and $\eta\ne\alpha$.
The number $t=|F|$ is exactly
\[
 t=\begin{cases}
 1,&\eta=1,\\
 2,&\eta^2=1,\ \eta\ne1,\\
 3,&\eta^2\ne1.
 \end{cases}                                               \tag{B5}
\]
All three labels in (A3) remain even when their images coincide.
\end{theorem}
\begin{proof}
$W$ is the full inverse image of $S$ inside $\Gamma$, so membership of
any target is equivalent to membership of its coset. Its involution
$\alpha$ cannot be one because $K\subset W$. If $\eta=\alpha$, an exterior
target is one, also impossible. Equality of an exterior class with the middle
class is equivalent to $\eta=1$. Equality of the two exterior classes is
equivalent to $\eta^2=1$. These prove every case of (B5).
\end{proof}
For a prime residual $R\equiv3\pmod4$, a joint-failure stabilizer has odd
order, since every even subgroup of the cyclic unit group contains $-1$.
Thus every odd \emph{ambient} index is already rescued by M. In particular,
the earlier one-target cubic defect is not a joint obstruction
\cite[\S\S4--8]{BryanTwo}. For composite $R$, the involution collision in
the second row of (B5) is real and cannot be removed by that cyclic argument.

\section{Bounds at every original factor cut}
Call $i$ active when $g_i\ne1$, and put
\[
 E_K=\sum_{i\text{ active}}e_i.                            \tag{C1}
\]
An empty shell has the following unconditional structural restrictions.

\begin{theorem}[Joint and cut budgets]\label{thm:cuts}
For every active index,
\[
 \operatorname{ord}(g_i)>2e_i+1,
 \qquad
 \boxed{2E_K\le d-1-|F|.}                               \tag{C2}
\]
For a nonempty subset $J$ of the active indices define
\[
 C_J=\langle g_i:i\in J\rangle,\quad
 S_J=\prod_{i\in J}\{g_i^{-e_i},\ldots,g_i^{e_i}\},\quad
 U_J=\prod_{i\notin J}\{g_i^{-e_i},\ldots,g_i^{e_i}\},
\]
\[
 F_J=(F U_J^{-1})\cap C_J.
\]
Then $1\notin F_J$, $S_J\cap F_J=\varnothing$, and
\[
 \boxed{2\sum_{i\in J}e_i
 \le |C_J|-1-\max\{1,|F_J|\}.}                          \tag{C3}
\]
In particular, for every nontrivial cyclic subgroup $C\le A$,
\[
 \boxed{2\sum_{\langle g_i\rangle=C}e_i\le |C|-2.}       \tag{C4}
\]
\end{theorem}
\begin{proof}
Any period of a factor subproduct remains a period after multiplying the
other factors. Since $S$ is aperiodic, so is every such subproduct.
If $2e_i+1\ge\operatorname{ord}(g_i)>1$, its interval is the entire cyclic
subgroup, a contradiction. Translate each interval to $[0,2e_i]g_i$ and
adjoin its $2e_i$ two-point factors $\{1,g_i\}$. Every step grows the
current set strictly: otherwise $g_i$ is a period, which would persist to
$S$. Thus $|S|\ge1+2E_K$. Together with $|S|\le d-|F|$ this proves (C2).

The same growth proof gives $|S_J|\ge1+2\sum_J e_i$. Joint failure says
exactly that $S_J$ avoids every inverse translate of a target by an actual
remaining word; it therefore avoids $F_J$. Moreover $S_J$ cannot equal the
nontrivial group $C_J$, since that would create a period of $S$.
It misses at least $\max\{1,|F_J|\}$ elements of $C_J$, proving (C3).
For (C4), group the indicated actual factors, whose generated group is $C$,
and apply the same at-least-one-hole argument.
\end{proof}
These proofs specialize the strict-growth mechanism underlying Kneser's
addition theorem \cite[Theorem 1]{DeVos}; the earlier single-target and
prime-nonresidue versions are attributed to \cite{BryanDefect,BryanTwo}.
No use is made of a density statement or a global conclusion of any source.
The complement $U_J$ uses the actual remaining bounded factors; independently
chosen residues at different prime powers of $R$ do not replace it.

For fixed $d$, (C2) leaves at most $\lfloor(d-2)/2\rfloor$ active prime
occurrences. All prime factors inside $K$, their exponents, and their original
coefficient distribution remain unrestricted data. This is a finite reduction
of the active quotient profile, not a uniform bound on $d$ for every shell.

\section{Complete small-effective-index normal forms}
Let
\[
 B=\prod_{q_i\in K}q_i^{e_i},\qquad a=B\,a_{\rm act}.
                                                               \tag{D1}
\]
The definition is from the actual period subgroup, not an assumed quadratic
or power-residue subgroup. In the following tables powers of $g$ are quotient
classes; actual prime labels and inverse orientations are retained.

\begin{theorem}[Classification through order eight]\label{thm:low}
If an original shell fails jointly and $d\le8$, then
\[
 \boxed{W(B)=K,\qquad \Omega(a_{\rm act})\le2.}           \tag{D2}
\]
The entire list of pure cases $a_{\rm act}=1$ is
\[
\begin{array}{c|c|l}
 d&A&\text{coefficient data}\\\hline
 2&C_2&\eta=1\\
 4&C_2\times C_2&\alpha,z\text{ independent},\ \eta=1\\
 6&C_6&\operatorname{ord}(z)\in\{3,6\},\ \operatorname{ord}(\eta)=3\\
 8&C_8&\operatorname{ord}(z)=8,\ \alpha=z^4,\ \eta=z^2\\
 8&C_4\times C_2&\operatorname{ord}(z)=4,\ \alpha\notin\langle z\rangle,\ \eta=z^2.
\end{array}                                                   \tag{D3}
\]
The entire list of nonpure cases is
\[
\begin{array}{c|c|c|c}
 A&\text{active occurrences}&S&\text{phase requirement}\\\hline
 C_6&1,\ \operatorname{ord}(g)=6&\{1,g^{\pm1}\}&\eta=g^{\pm1}\\
 C_6&2,\text{ all in }\{g,g^{-1}\},\ \operatorname{ord}(g)=6
     &\{1,g^{\pm1},g^{\pm2}\}&\eta=1\\
 C_8&1,\ \operatorname{ord}(g)=4&\{1,g^{\pm1}\}&\eta=1\\
 C_8&1,\ \operatorname{ord}(g)=8&\{1,g^{\pm1}\}&\eta=g^{\pm1}\\
 C_8&2,\text{ all in }\{g,g^{-1}\},\ \operatorname{ord}(g)=8
     &\{1,g^{\pm1},g^{\pm2}\}&\eta=1\\
 C_4\times C_2&1,\ \operatorname{ord}(g)=4&\{1,g^{\pm1}\}
     &\alpha\notin\langle g\rangle,\ \eta=g^{\pm1}.
\end{array}                                                   \tag{D4}
\]
Two occurrences mean one actual prime of exponent two or two distinct primes
of exponent one. Relation (B3) is imposed throughout, and the quotient is
generated by $\alpha,z,g_i$. Conversely, these tables together with $W(B)=K$
produce joint failure and make $K$ the actual stabilizer.
\end{theorem}
\begin{proof}
Failure makes $d$ even. Active elements have order at least four by (C2).
At $d=2$ there are none. At $d=4$, only a cyclic order-four generator could
be active, once. Its support is $\{1,g,g^{-1}\}$, missing only
$\alpha=g^2$. Joint failure would force $\eta=1$, but (B3) makes
$\eta=z^2g^{\pm1}$ an odd power of $g$, impossible. Thus both cases are pure.
A pure quotient is generated by $\alpha,z$, with $\eta=z^2$.
At order four it cannot be cyclic, because a generator $z$ would give
$\eta=\alpha$ and an exterior hit. This gives the first two rows of (D3).

At order six, active elements cannot have order two or three. The budget
allows at most two occurrences, all in the unique inverse pair of generators.
One occurrence gives $S=\{1,g,g^{-1}\}$. By (B3), $\eta$ is an odd power
of $g$; the value $g^3=\alpha$ produces the identity target and is forbidden.
This leaves $\eta=g^{\pm1}$. Two occurrences give the five-element interval,
whose only missing element is $\alpha$; both exterior targets can miss only
when $\eta=1$. A pure order-six quotient has $z$ of order three or six,
otherwise $\alpha,z$ would generate a smaller group. This proves the order-six rows.

At order eight the possibilities are $C_8$, $C_4\times C_2$, and $C_2^3$.
The last has no active elements and cannot be generated purely by two elements.
In $C_4\times C_2$, every active element has order four and capacity one.
Two from the same cyclic subgroup already fill it, a forbidden period.
Two from its two different cyclic subgroups have seven-element sumset.
A joint miss would then require $\eta=1$. Their central product lies outside
the square subgroup $\{1,h^2\}$, whereas $z^2$ lies inside it, contradicting
(B3). Three active occurrences repeat a cyclic subgroup. For one occurrence,
$\eta=z^2g$ is $g$ or $g^{-1}$. If $\alpha\in\langle g\rangle$, an exterior
target belongs to $S$; if it is outside, all three targets are outside $S$.
The pure case gives the last row of (D3).

In $C_8$, an order-four active element has capacity one. Mixed order-four
and order-eight elements give seven support elements but an odd central
product, hence an odd $\eta$; joint failure of a seven-element support would
require $\eta=1$. Two order-eight directions outside one inverse pair fill
$C_8$. Three aligned occurrences give seven elements but again odd $\eta$.
Two order-four occurrences fill their subgroup. Only the listed cases remain.
For a sole order-four element, testing $\alpha\eta^{\pm1}$ forces $\eta=1$.
For a sole generator, it forces $\eta=g^{\pm1}$. For two aligned generator
occurrences, it forces $\eta=1$. The pure cyclic case requires $z$ itself
to generate $C_8$, proving the remaining rows.

It remains to prove (D2), not assume inactive saturation. In each active
case at least one extreme quotient exponent has exactly one preimage in the
active exponent box. The full fine support in that coset is a translate of
$W(B)$. Since it is a nonempty complete $K$-coset by actual periodicity,
$W(B)=K$. In a pure case $W\subseteq K$ and $K\subseteq W$ give the same
conclusion. Finally, every tabulated $S$ is aperiodic and avoids its stated
three labelled targets. If $W(B)=K$, its full inverse image is exactly $W$;
this proves the converse and the actual-stabilizer assertion.
\end{proof}

\begin{corollary}[The first active obstruction and the collision threshold]
If $W$ is not a subgroup and the shell fails jointly, then $d\ge6$.
If in addition $|F|=2$, then $d\ge12$.
\end{corollary}
\begin{proof}
The first claim follows from the table at orders two and four. No nonpure
row at orders six or eight has two distinct targets. At order ten the abelian
group is cyclic, so its only nontrivial involution is $\alpha$; (B5) would
force $\eta=\alpha$, which gives a hit. Odd orders cannot occur.
\end{proof}

The multiplicities in (D4) are also exact. Put $T=\tau(B^2)$.
For one simple active prime, the three quotient masses are $(T,T,T)$.
For one active square prime they are $T(1,1,1,1,1)$.
For two distinct active simple primes they are $T(1,2,3,2,1)$, in
ascending exponent order $-2,-1,0,1,2$. The fine coefficients remain
translated sums of the full function $\mu_B$; they are not obtained by
uniformly distributing those totals over $K$.

\section{The next-turn prime-shift input, with its precise scope}
Suppose now that $R=q$ is prime, $q\equiv3\pmod4$, $q<p$, and
$(p/q)=-1$. Then $a\equiv p/4$ is a nonresidue modulo $q$.
The group $\Gamma$ is cyclic with exactly one factor of two in its order;
$K$ has odd order, hence is contained in the quadratic residues. At least
one original prime factor is therefore active. Moreover $\eta$ is a
nonresidue in the quotient while $z^2$ is a square.

The order-six case of Theorem~\ref{thm:low} consequently forces
\[
 \boxed{
 a=Br,\quad r\text{ prime},\quad v_r(a)=1,\quad W(B)=K,\quad
 \operatorname{ord}(rK)=6,\quad pK=rK\text{ or }r^{-1}K.}  \tag{E1}
\]
All prime factors of $B$ lie in $K$. The active prime is the unique nonresidue
factor modulo $q$. For an odd prime $s\mid a$, the equation $4a=p+q$
gives $p\equiv-q\pmod s$. Reciprocity, with $p\equiv1\pmod4$ and
$q\equiv3\pmod4$, gives
\[
 (s/p)=(p/s)=(-1/s)(q/s)=(s/q).
\]
For $s=2$, its occurrence in $a$ gives $q\equiv-p\pmod8$, and the same
identity follows from the supplementary law. Thus the active prime is also
the unique nonresidue factor modulo $p$. The next factor edge is determined,
with its coefficient block retained. If the \emph{ambient} index is six, $K=G^6$; for an effective
index-six quotient inside a proper $\Gamma$, that equality with $G^6$ must
not be asserted. The ambient prime case is inherited from
\cite[\S\S4--8]{BryanSix}; the exact inactive block and fine fibres are part
of the present reconstruction.

The two phases in (E1) are respectively
\[
 4\in K,\qquad 4r^2\in K.                                \tag{E2}
\]
If $p\equiv1\pmod8$ and $q\equiv7\pmod8$, then $a$ is even. The prime
$2$ is a quadratic residue modulo $q$ and cannot be the exceptional $r$;
therefore $2\mid B$, $2\in K$, and only the first phase is possible.
For $q\equiv3\pmod8$, $a$ is odd and both alternatives require their own
check. No theorem that all such prime shifts cannot fail simultaneously
is claimed. The existing source even records consecutive primitive defects
\cite[\S\S1--3]{BryanChain}; no contradiction from adjacency alone is used.

\section{Exact extension coordinates and the retained coefficient kernel}
Choose a section $\sigma:A\to\Gamma$ with $\sigma(1)=1$, for example the
least positive representative. Each original unit has the unique coordinates
\[
 g=\sigma(c)k,\qquad c=\pi(g),\quad k=g\sigma(c)^{-1}\in K.
\]
Define
\[
 \omega(c,d)=\sigma(c)\sigma(d)\sigma(cd)^{-1}\in K.
\]
The exact multiplication is
\[
 \boxed{(c,k)(d,l)=(cd,kl\omega(c,d)).}                    \tag{F1}
\]
Expanding the representatives proves
$\omega(c,d)\omega(cd,e)=\omega(d,e)\omega(c,de)$ and
$\omega(1,c)=\omega(c,1)=1$. The displayed product therefore retains every
coefficient and has inverse $g\mapsto(\pi(g),g/\sigma\pi(g))$.
Replacing $\sigma$ by $\sigma(c)\theta(c)$ changes the kernel coordinate to
$k\theta(c)^{-1}$ and the cocycle by
$\theta(c)\theta(d)\theta(cd)^{-1}$. It is a coordinate change, not a
claim that the extension splits.

The complete fine count is
\[
 \nu(c,k)=\mu(\sigma(c)k),\qquad
 \bar\mu(c)=\sum_{k\in K}\nu(c,k).                       \tag{F2}
\]
The integer quotient kernel is generated by
$e_{\sigma(c)k}-e_{\sigma(c)}$ for $k\ne1$. Given its quotient value and
all nonreference coefficients, this star-difference expansion reconstructs
the original vector uniquely. On the original exponent source, the analogous
kernel has one difference per nonreference $\beta$ in its actual residue
fibre. No representative is introduced outside $\mathcal B$.

Actual support periodicity makes Boolean target membership equivalent in the
quotient, but it does not make (F2) a uniform distribution. To return an exact
fine target, retain its kernel coordinate, and use the original suffix count:
a positive convolution coefficient has some available next exponent whose
suffix coefficient is positive. Iterating reconstructs $\beta$, then (A2)
and (A5). All choices give the full finite fibre. Quotienting before retaining
these data loses the fine witness and its multiplicity.

In the original SplitZero convention $G(A)=\{\tau\}\sqcup A^\bullet$,
a nonempty source fibre with zero image is supported zero, not $\tau$.
The quotient and its incoming star relations above preserve that distinction.
The present finite extension cocycle is not identified with an analytic
Gamma or theta norm, or with a source-wide formal certificate \cite{SplitZero}.

\section{Original arithmetic certificates and the nonclaims}
At $p=3361,a=848=2^4\cdot53,R=31$,
\[
 K=\{1,2,4,8,16\},\quad |\Gamma/K|=6,\quad a_{\rm act}=53.
\]
Both channels are empty. The inactive box has coefficient one at $1$ and
coefficient two at each other member of $K$. Each of the three occupied
quotient cosets has mass nine, not five equal coefficients $9/5$.
Here $pK=53K$ and $2\in K$.

At $p=5569,a=1397=11\cdot127,R=19$, the other phase occurs:
\[
 K=\{1,7,11\},\quad pK=127^{-1}K,\quad \Gamma=C_{18}.
\]
For the section $\sigma(j)=2^j$, $0\le j<6$, the actual carry is
\[
 \boxed{\omega(i,j)=7^{\lfloor(i+j)/6\rfloor}\pmod{19}.} \tag{F3}
\]
Every lift of the quotient generator has order eighteen, so no homomorphic
section exists. A direct product $C_6\times C_3$ would erase a real cubic
carry. All nine original words and all three target fibres are recorded.

The separate exterior defect at the hard prime $p=22129$ has
$a=5537=7^2\cdot113,R=19$ and $W=K$ of ambient index three. E is empty,
but the two original M states include
\[
 (u,h,r_0,s_0,\lambda)=(49,49,1,113,6),\qquad
 \frac4{22129}=\frac1{5537}+\frac1{735169638}+\frac1{6505926}.
\]
Thus this one-channel cubic defect is not an unresolved joint case.

The composite-residual collision $p=109,a=31,R=15$ has $K=\{1\}$,
$\eta=4$, and targets $14,11,11$: exactly two classes, both channels empty.
The distinct effective/ambient example $p=6340273,a=1259^2,R=51$ has
\[
 W=K=\{1,35\},\quad \mu(1)=3,\quad\mu(35)=2,
 \quad |\Gamma/K|=8,\quad |G/K|=16.
\]
It realizes the pure cyclic order-eight case. Cyclic prime-residual reasoning
cannot be transferred to this composite modulus.

The complete source-hashed replay covers all $73\,798$ first-half shells at
$211$ primes $p\equiv1\pmod4$, $p\le3000$, and $1\,877\,004$ original divisor
vectors, with both canonical channel tests on each. It records $5700$ exterior
and $5268$ middle states in $3918$ occupied shells. All $69\,880$ jointly
empty shells undergo the actual-stabilizer and every active-factor-cut test;
all $459$ such shells of effective index at most eight satisfy the normal forms
and inactive saturation. Abstract tests exhaust $596$ aperiodic,
coefficient-anchored profiles of capacity at most three, including $146$
failures, in every abelian group of order $2,4,6,8$. The capacity bound follows
from (C2). The main verifier makes $2\,153\,122$ named \texttt{check()} calls,
in addition to eight successful negative-control rejections and other explicit
precondition tests. A separate direct-divisor implementation imports neither
it nor the antecedent software and recomputes the arithmetic of every supplied
ledger row. It does not independently enumerate the ledger's prime/shell
domain, so it is not a second completeness certificate for that row set.
Execution hashes, the exact separate check count, and normal/optimized output
identities are in the receipt. These finite checks neither prove the structural
theorem by sampling nor constitute independent mathematical review. No new
prime-coverage range is asserted.

The required exit of Turn 1 is the unconditional joint reduction (B)--(D),
its complete low-index normal forms, and its integer coefficient return.
Turn 2 must connect this exact object to the original shifted-factor graph.
In particular, primitive sextic failures exist; the next theorem cannot
assert that an individual such shell is impossible. No uniform upper bound
on the effective index, global contradiction, historical-priority claim,
independent mathematical review, or new Lean build is asserted.
